% ============================================================
% CSMATH 2026 课程论文 · 导言区
% 浙江大学计算机科学与技术学院 · 计算机应用数学
% ============================================================

% --- 基础宏包 ---
\usepackage{xcolor}

% --- 数学宏包 ---
\usepackage{amsmath,amssymb,amsthm}
\usepackage{mathtools}

% --- 图表与图形 ---
\usepackage{graphicx}
\usepackage{booktabs}
\usepackage{caption}
\usepackage{subcaption}

% --- 算法 ---
\usepackage[ruled,linesnumbered]{algorithm2e}
\SetKwComment{Comment}{$\triangleright$\ }{}

% --- 定理环境 ---
\newtheorem{theorem}{定理}[section]
\newtheorem{lemma}[theorem]{引理}
\newtheorem{proposition}[theorem]{命题}
\newtheorem{corollary}[theorem]{推论}
\newtheorem{definition}{定义}[section]
\newtheorem{remark}{注记}[section]
\newtheorem{property}{性质}[section]
\theoremstyle{definition}
\newtheorem{example}{例}[section]

% --- 常用数学运算符 ---
\DeclareMathOperator*{\argmin}{arg\,min}
\DeclareMathOperator*{\argmax}{arg\,max}
\DeclareMathOperator*{\E}{\mathbb{E}}
\DeclareMathOperator*{\Var}{Var}
\DeclareMathOperator*{\Cov}{Cov}
\DeclareMathOperator{\diag}{diag}
\DeclareMathOperator{\tr}{tr}
\DeclareMathOperator{\rank}{rank}

% --- 常用符号 ---
\newcommand{\R}{\mathbb{R}}
\newcommand{\C}{\mathbb{C}}
\newcommand{\N}{\mathbb{N}}
\newcommand{\Z}{\mathbb{Z}}
\newcommand{\X}{\mathcal{X}}
\newcommand{\Y}{\mathcal{Y}}
\newcommand{\F}{\mathcal{F}}
\newcommand{\Loss}{\mathcal{L}}
\newcommand{\norm}[1]{\left\lVert#1\right\rVert}
\newcommand{\inner}[2]{\left\langle#1, #2\right\rangle}
\newcommand{\grad}{\nabla}

% --- 行内标注（写作过程中使用，提交前禁用）---
\newcommand{\red}[1]{{\color{red}#1}}
\newcommand{\todo}[1]{{\color{red}#1}}
\newcommand{\TODO}[1]{\textbf{\color{red}[TODO: #1]}}
% 提交版本取消注释下面两行以隐藏 TODO
% \renewcommand{\TODO}[1]{}
% \renewcommand{\todo}[1]{#1}

% --- 微排版优化（定稿前取消注释）---
% \usepackage{microtype}

% --- 其他 ---
\usepackage{enumitem}
\usepackage{multirow}
\usepackage{url}
